\section{Formulation for LaMET and Hybrid Renormalization}
\label{sec:analytic}

The leading-twist gluon PDFs $g(x)$ of a hadron is defined as~\cite{Collins:1989gx},
%
\begin{align}\label{eq:LCPDF}
g(x,\mu) = \frac{1}{ 2 x P^+} \int_{-\infty}^\infty \frac{\mathrm{d} \xi^-}{2\pi} \,& e^{-i x P^+ \xi^-} \langle P | F_a^{+\mu}(\xi^-) \notag\\
&
 \times \mathcal{W}_{ab}(\xi^-, 0) \, F_{b,\mu}^+(0) | P \rangle,
\end{align}
where $|P\rangle$ denotes unpolarized hadron with momentum $P$ along the $z$-direction.  Here, $x$ is the longitudinal momentum fraction carried by the gluon, and $\mu$ denotes the renormalization scale in the $\overline{\rm MS}$ scheme. The light-front coordinates are given by $\xi^\pm=(\xi^t\pm \xi^z)/\sqrt{2}$ and gauge invariance is ensured by the Wilson line,
\begin{align}
\mathcal{W}_{ab}(\xi^-, 0)= \mathcal{P}\,{\rm exp}\big[ -i g \int_0^{\xi^-} \mathrm{d} \eta^- \, A_c^+(\eta^-) t^c \big]_{ab},
\end{align}
where the indices $\{a,b\}$ denote the color components in the adjoint representation of SU(3), and $t^c$ are the generators of the gauge group.

According to LaMET, the gluon PDF can be extracted by calculating a spatially correlated matrix element:
\begin{align}
h^{\text{B}}(z,P_z)= \langle P | &F^{ti}(z)\mathcal{W}(z,0)F^t_{i}(0) \notag\\
&-F^{ij}(z)\mathcal{W}(z,0)F_{ij}(0)| P \rangle,
\end{align}
where $i$ and $j$ run over the transverse indices $\{x,y\}$, only.
Hereafter, we adopt the convention of omitting color indices for notational simplicity.

Although the local case does not strictly admit a multiplicative renormalization, we do not incorporate information from the local region, as the result is normalized and therefore insensitive to local contributions. The bare non-local quasi-light-front (LF) correlation can be multiplicatively renormalized as~\cite{Zhang:2018diq}
%
\begin{align}
h^{\text{B}}(z,P_z)=Z_L e^{-\bar{\delta} m |z|} \,h^{\text{R}}(z,P_z).
\end{align}
where $h^{\text{R}}(z,P_z)$ is the renormalized quasi-LF correlation that contains both $z$-independent logarithmic UV divergences, captured by $Z_L$, and $z$-dependent linear divergences from Wilson-line self-energy, denoted by $\bar{\delta} m$.
To remove these divergences, we adopt the hybrid renormalization scheme~\cite{Ji:2020brr}, treating short and long distances separately:
at short distances, we apply the ratio scheme~\cite{Radyushkin:2018cvn} using rest-frame matrix elements.
While at long distances, we use mass renormalization by fitting the zero-momentum matrix $h^{\text{B}}(z,0)$ to capture linear divergence effects.
The resulting renormalized quasi-LF correlation is expressed as:
\begin{equation}\label{eq:hybridscheme}
h^{\mathrm{R}}(z, P_z) =
\begin{cases}
\!\displaystyle\!
\frac{h^{\mathrm{B}}(0,0)}{h^{\mathrm{B}}(0,P_z)} \frac{h^{\mathrm{B}}(z,P_z)}{h^{\mathrm{B}}(z,0)} & z \leq z_s \\[12pt]
\!\displaystyle\!
\frac{h^{\mathrm{B}}(0,0)}{h^{\mathrm{B}}(0,P_z)}  \frac{h^{\mathrm{B}}(z,P_z)}{h^{\mathrm{B}}(z_s,0)} e^{(\delta m + m_0)(z - z_s)} & z > z_s
\end{cases}
\end{equation}
where $z_s$ separates short and long distances, the terms $\delta m$ and $m_0$ correspond to the mass renormalization and the renormalon ambiguity, respectively, and are introduced to cancel the linear divergence arising from the self-energy of the Wilson line. In the limit $z_s \to \infty$, the expression reduces to the standard ratio-scheme renormalization, which neglects the Wilson line self-energy contribution.
\UPDATE{Using the double ratio in the short distance is a good method for handling discretization errors associated with the non-commutativity of the $z \rightarrow 0$ and lattice spacing $a \rightarrow 0$ limits~\cite{Ji:2020brr}, though this scheme will cause the resulting gluon PDF to be normalized by its first moment $\langle x\rangle_g$.
We can see in a previous study of the continuum limit of the gluon PDF that discretization errors in the short distance using the ratio scheme are well under control within current statistical errors~\cite{Fan:2022kcb}.
It has been suggested that more control over discretization effects can be achieved using the self-renormalization method~\cite{LatticePartonLPC:2021gpi}, though this requires the use of data from multiple lattice spacings.}

\begin{figure*}[htbp]
\centering
        \includegraphics[width=.48\textwidth]{images/W3_m0vz.pdf}
\centering
        \includegraphics[width=.48\textwidth]{images/H5_m0vz.pdf}

    \caption{Fit results for the $\delta m + m_0$ at $\mu = 2 $ GeV for varying central point $z$ for each hadron for the two smearings, Wilson-3 (left) and HYP5 (right).
    The error bands give statistical fluctuations only.
    The points represent the choice of $\delta m + m_0$ that is used in the subsequent hybrid renormalization.}
    \label{fig:m0_vs_z}
\end{figure*}

In particular, the availability of multiple lattice spacings allows for a separate determination of $\delta m$ and $m_0$ via fitting procedures~\cite{Ji:2020brr, LatticePartonLPC:2021gpi, Zhang:2023bxs}.
However, in the case of a single lattice spacing, it is more practical to fit only their sum, $\delta m + m_0$, by matching the $P_z = 0$ bare matrix elements to the perturbatively calculated ``Wilson coefficients". The Wilson coefficient can be derived from Ref.~\cite{Balitsky:2021qsr}, gives
\begin{align}\label{eq:wil-coef}
\mathcal{H}(z,\mu) = 1 + \frac{\alpha_s C_A}{2\pi} \left( \frac{5}{6} \ln \left( \frac{z^2 \mu^2 e^{2\gamma_E}}{4} \right) + \frac{3}{2} \right).
\end{align}
Then we determine the combined parameter $\delta m + m_0$ in the short-distance region by fitting to the relation:
\begin{equation}\label{eq:dm_fit}
(\delta m + m_0) z  + I_0 \approx \ln\left[ \frac{\mathcal{H}(z,\mu)}{h^{\mathrm{B}}(z,0)} \right],
\end{equation}
where $I_0$ denotes a $z$-independent constant.

The renormalized quasi-LF correlation function can be transformed into the so-called quasi-PDF via a Fourier transform. The gluon quasi-PDF is related to the light-front PDF through a factorization formula as following:
\begin{align}\label{eq:mtcheq}
 x\tilde{g}\left(x, P_z\right)= & \int_{-1}^1 \frac{d y}{|y|} C\left(\frac{x}{y}, \frac{\mu}{y P_z}\right)  y g(y, \mu)  +h.t. ,
\end{align}
where $g(y, \mu)$ is the light-front gluon PDF at the renormalization scale $\mu$, $C\left(x/y, \mu/(y P_z)\right)$ represents the perturbative matching coefficient, and $h.t.$ denotes higher-twist corrections suppressed by powers of $1/P_z$.
\UPDATE{Notice that we neglect the contribution of mixing with quarks in this work, as it has been shown to have an effect smaller than $\sim10$\% in the pseudo-PDF studies~\cite{Fan:2022kcb}, which is well within statistical precision and smaller than the lattice systematic effects that we explore in this paper.}
The perturbative matching coefficient in coordinate space up to the NLO in the ratio scheme can be found in Refs.~\cite{Balitsky:2019krf, Balitsky:2021qsr}. Then we obtain the corresponding momentum-space kernel directly via a Fourier transform, following the procedure outlined in Ref.~\cite{Yao:2022vtp}. The matching coefficients in the ratio and hybrid schemes are given as follows:

\begin{equation}\label{eq:kernel}
C^{\text{ratio}}\left(\xi, \frac{\mu}{p_z}\right)
= \delta(1-\xi)
+ \frac{\alpha_s C_A}{2\pi}
\left\{
\begin{array}{ll}
\left(
2 \frac{(1 - \xi + \xi^2)^2}{1 - \xi} \ln\left(\frac{\xi}{\xi-1}\right)
+\frac{11 - 28\xi + 18\xi^2 - 12\xi^3}{6(1-\xi)}
\right)_{+(1)}, & \xi > 1 \\[1.5ex]
\left(
2 \frac{(1 - \xi + \xi^2)^2}{1 - \xi} \left[ -\ln(\frac{\mu^2}{4{p_z}^2}) + \ln\left[\xi(1-\xi)\right] \right]
- \frac{15 - 56\xi + 102\xi^2 - 96\xi^3 + 48\xi^4}{6(1-\xi)}
 \right)_{+(1)}, & 0 < \xi < 1 \\[1.5ex]
\left(
-2 \frac{(1 - \xi + \xi^2)^2}{1 - \xi} \ln\left(\frac{-\xi}{1-\xi}\right)
- \frac{11 - 28\xi + 18\xi^2 - 12\xi^3}{6(1-\xi)} \right)_{+(1)}, & \xi < 0
\end{array}
\right.
\end{equation}

\begin{equation}
C^{\text{hyb.}}\left(\xi, \frac{\mu}{p_z}\right) =C^{\text{ratio}}\left(\xi, \frac{\mu}{p_z}\right) + \frac{\alpha_s C_A}{2\pi} \frac 5 6 \bigg[-\frac{1}{|1-\xi|} +\frac{2\text{Si}((1-\xi)z_s p_z)}{\pi(1-\xi)}\bigg]_+ .
\end{equation}

\noindent where $\xi = \frac{x}{y}$ and $p_z$ denotes the momentum carried by the parton, which is given by $y P_z$.
